%% GENERATED by tools/gen_dead_ends.py from tools/dead_ends.tsv — do not edit by hand.
\renewcommand{\arraystretch}{1.12}\small
\begin{longtable}{@{}p{0.55cm}p{3.1cm}p{4.35cm}p{5.4cm}@{}}
\toprule \textbf{\#} & \textbf{Dead end} & \textbf{Approach} & \textbf{Why it fails} \\ \midrule \endfirsthead
\toprule \textbf{\#} & \textbf{Dead end} & \textbf{Approach} & \textbf{Why it fails} \\ \midrule \endhead
\midrule \multicolumn{4}{r}{\emph{continued on next page}} \\ \endfoot
\bottomrule \endlastfoot
\multicolumn{4}{@{}p{13.4cm}@{}}{\textsc{Orbit specificity} (OS, 10 entries) --- \emph{population or ensemble statistics do not determine what one orbit does}.} \\[2pt]
9 & Chebotarev / Kummer route {\footnotesize(S10)} & Kummer extensions plus Chebotarev density (Booker--Simon style) to prove DH algebraically without character sums & Chebotarev not in Mathlib ($\sim$5000 lines); even if available it gives population density (SieveEquidistribution), not orbit-specific SieveTransfer. Abelian case degenerates to Dirichlet in APs (Session 298) \\
18 & Target-diversity heuristic (SieveTransfer) {\footnotesize(S13)} & Heuristic ``each step stays in $\ker\chi$ with probability $1/d$'' to bound kernel-run lengths & Heuristic only: needs an Alladi-type theorem for the specific EM products; reduces to proving $\mathrm{prod}(n)+1$ generic w.r.t. small factors = SieveTransfer. Superseded: GenericLPFEquidist false, SieveTransfer vacuous (\#160) \\
30 & Information-theoretic argument for BRE {\footnotesize(S16)} & Kolmogorov-complexity / counting arguments for the BRE gap & BRE failure has ensemble probability $(s/d)^N$ but EM is one deterministic sequence (zero entropy); ensemble-to-specific gap; restates decorrelation without content \\
32 & Bombieri--Vinogradov for EM subsequence {\footnotesize(S17)} & Use BV (average over moduli) for the multiplier primes & BV applies to all primes in APs, not to the greedy EM subsequence; also Mathlib-blocked (\#55--57); EM bridge \code{BVImpliesMMCSB} open, MultCSB $\not\Rightarrow$ MMCSB (\#58) \\
90 & Population statistics do not determine orbit {\footnotesize(S59)} {\footnotesize[revival 3]} & Derive CME (later DSL, UCE $\to$ CME, Borel--Cantelli, CRT spreading, MinFacUnbiased, sieve orbit control, Deligne/Haar equidistribution) from population-level results on the least prime factor. & Four-Layer Gap: no technique spans population $\to$ individual, unconditional $\to$ conditional, static $\to$ growing, counting $\to$ distribution. Witness: periods $(2,2,3,3)$ and $(2,3,2,3)$ in $(\mathbb Z/5)^\times$ have identical statistics; only the first hits $-1$. Superseded in part: PE/UCE false (\#160), bridges vacuous. \newline {\footnotesize\color{leanink}$\checkmark$\,\code{OrbitBarrier.population\_does\_not\_determine\_hitting}} \\
127 & Weil bound does not close FF-DSL {\footnotesize(S166)} {\footnotesize[revival 2]} & Euclid--Mullin over $\mathbb{F}_p[t]$, where the Weil RH is a theorem, hoping unconditional ANT closes the function-field DSL. & Weil gives population equidistribution of irreducibles only; $\mathrm{ffProd}(n)+1$ is one specific polynomial, the walk sum is a sum of products of character values (not a Weil sum), and the relevant population is $O(1)$; orbit specificity (\#90) and the Four-Way Blocker are unchanged. \\
148 & Consumption ledger is one-sided {\footnotesize(S299)} {\footnotesize[revival 1]} & Shield-ledger incidence module: count small-prime consumption by an avoiding orbit and contradict a budget. & The ledger only yields caps (\code{hittingSet\_ncard\_le\_appearing}, starvation); a contradiction needs a lower bound on spending, which requires orbit control, exactly what the consumption discipline forbids. Detection strength is zero (tail class $0$ unconditionally via \code{exists\_tail\_coprime}); the Gap is false, not hard. \newline {\footnotesize\color{leanink}$\checkmark$\,\code{hittingSet\_ncard\_le}} \\
158 & Tail-identity Borel--Cantelli is \#90 again {\footnotesize(S307)} {\footnotesize[revival 0]} & \code{FourPointPCVImpliesDSL}: the proved tail identity makes $\prod M$ an ensemble member; a density-$O(1/K^2)$ bad-set bound plus Borel--Cantelli should give cofinally good windows for the orbit. & \code{standard\_tail\_not\_bad} needs the bad set to have $\mathrm{card}=0$, not density $\to 0$; a density-zero bad set may contain the sparse sequence $\{\prod M\}$ at no cost. The identity turns an orbit statement into an ensemble-member statement, but specificity reappears as which members are good (\#58/\#117). \\
172 & Sure compensator as evidence of coverage progress {\footnotesize(S313)} {\footnotesize[revival 0]} & Read \code{pathwise\_compensator} ($\sum_{k<n} S_k \ge (c_1/2)n$, sure and per-path) as forcing the orbit to make progress toward covering the residues mod $q$. & Wrong sign. $S_k = \prod_r (1-\rho_r)$ is \emph{maximal} ($\equiv 1$) exactly when no position is ever new --- i.e. for the orbits that capture nothing. Per-path, ``large $S_k$'' means ``no progress toward coverage''; the opposite (ensemble) reading requires the seed residue to be uniform in its box, which is \#90. Both proved sure bounds --- charge $\le$, survival $\ge$ --- constrain the half-line \emph{opposite} to capture-freeness: the configurations satisfying the whole layer are downward-closed in charge, and the capture-free ones sit at the extreme point (zero charge, $S_k\equiv 1$, boxes frozen). No upper bound on charge can exclude zero charge. \\
176 & The scale-uniform tail bound ``(N2)'' as the missing \S G input {\footnotesize(S317)} {\footnotesize[revival 1]} & Close the simultaneous-in-$q$ seed-average law by supplying, for every $\delta>0$, a $Q$ with $\#\{m\le X:\exists q>Q,\ q\nmid m,\ m \text{ never selects } q\}\le\delta X$ \emph{for all} $X$ (the input named (N2) in \code{docs/analysis/simultaneous\_in\_q\_scoping.md}, \S2.5), and combine it with \code{finite\_simultaneous\_density}. & Not a population statement. With $X=m$ and $\delta<1/m$ the bound leaves no room for the seed $m$ itself, so (N2) forces \emph{every} seed to miss only finitely many primes coprime to it (\code{GrowingRange.scaleUniformTail\_cofinite}); at $m=2$ it says the Euclid--Mullin sequence contains every sufficiently large prime (\code{scaleUniformTail\_cofinite\_mc}), far above the known floor MC $\Rightarrow$ RD $\Rightarrow$ (S) $\Rightarrow$ (C$\infty$). The honest input is (N2$'$), the same bound for $X\ge X_0(\delta)$; at scale $X$ it concerns the primes $q>X$, for which the seeds $m\le X$ are $X$ distinct points of $\mathbb{Z}/q$ and no period $M\le X$ sees the event --- orbit-specificity in the sense of \#90. What the population method does give is the growing-range form: some ineffective $Q\to\infty$ with a.a.\ seeds $m$ selecting every prime $q\le Q(m)$, $q\nmid m$ (\code{seed\_range\_never\_density}), the exact analogue of Tao's ``almost bounded values'' for Collatz. \newline {\footnotesize\color{leanink}$\checkmark$\,\code{scaleUniformTail\_cofinite\_mc}} \\
\midrule
\multicolumn{4}{@{}p{13.4cm}@{}}{\textsc{Decorrelation gap} (DG, 7 entries) --- \emph{multiplier-level cancellation or independence does not transfer to the walk}.} \\[2pt]
5 & Cofinal-cycle arguments plus SE {\footnotesize(S5)} & Derive a contradiction from cofinal-cycle constraints (positions, multipliers, death avoidance, product-one) plus SubgroupEscape; pumping $\ell > q-2$ & Cofinal orbit captures only marginal properties, all consistent with HH failure (§23 counterexamples); pumping gives $\ell \ge \Omega(q-1)+1$, logarithmic not linear. Marginal/Joint Barrier \\
58 & MultCSB does not imply MMCSB {\footnotesize(S28)} {\footnotesize[revival 0]} & Prove the abstract step \code{MultCSBImpliesMMCSB}: $o(N)$ multiplier character sums imply $o(N)$ walk character sums, splitting the BV route. & False: multipliers $(-1,-1,+1,+1)$ repeating give $S_M=O(1)$ but walk sum $N/2$; $(\omega,\omega^2,1)$ gives $\Theta(N)$. Order of multipliers matters for products, not sums; PED, NoLongRuns, DPED, Abel summation do not repair it. \newline {\footnotesize\color{leanink}$\checkmark$\,\code{OrbitBarrier.mult\_cancel\_not\_walk\_cancel}} \\
81 & CME does not imply HOD ($h\ge2$) {\footnotesize(S43)} {\footnotesize[revival 0]} & Route CME $\to$ HOD $\to$ (VdC) $\to$ CCSB using ConditionalMultiplierEquidist (multiplier character sums cancel on each walk-position fiber). & CME gives Dec ($h=1$) by fiber decomposition, but $h$-fold products couple consecutive multipliers through $w(n+1)=w(n)m(n)$; conditioning on a single position does not control these correlations. Hierarchy PED $<$ Dec $<$ CME $\le$ CCSB. \\
98 & CRT decorrelation gives no leverage {\footnotesize(S65)} & Use per-step CRT exogeneity (\code{crt\_multiplier\_invariance}: minFac$(P+1)$ depends only on $P$ mod primes $\ne q$), alias PBI, to get sequence-level decorrelation or DH. & Per-step exogeneity is not statistical independence: conditioning on $w(n)=c$ selects a biased time subsequence. CMI+FPE is CME. Counterexample: multipliers $\{2,3\}$ alternating in $(\mathbb Z/5)^\times$ generate the group but the walk cycles $\{1,2\}$, never hitting $-1$. \code{dh\_iff\_partial\_product\_surjective} false. \\
115 & Cofactor / accumulating-CRT joint barrier {\footnotesize(S109)} {\footnotesize[revival 0]} & View $P(n)$ in CRT space $\prod(\mathbb Z/p_i)$ hoping dimensional explosion yields independence mod $q$; later, the cofactor identity $w(n)+1=m(n)\,\mathrm{cofZ}(n)$ / CofactorEnsembleDecorrelation. & $P(n)$ is one locked point, not a generic one; the walk mod $q$ has dimension 1 (same $\{2,3\}$ in $(\mathbb Z/5)^\times$ counterexample). The cofactor is a joint quantity in bijection with the multiplier when alive, so CED is ensemble CME in other notation. \\
117 & MultCancel does not force WalkCancel {\footnotesize(S128)} {\footnotesize[revival 0]} & Prove MultCancelToWalkCancel for EM walks using squarefree accumulator, coprime cascade, minFac selection, CRT invariance and growth, completing the ensemble PT chain. & Multipliers alternating $\{2,3\}$ mod 5 with $\chi(2)=i$: multiplier sums vanish for even $K$, but $W_K=(K/2)(1+i)$, $|W_K|=\Theta(K)$. Compatible with every EM structural property; the transfer is equivalent to CCSB/CME. Sharpens \#58. \newline {\footnotesize\color{leanink}$\checkmark$\,\code{OrbitBarrier.mult\_cancel\_not\_walk\_cancel}} \\
123 & FourPointPCV: cross-time is not cross-modulus {\footnotesize(S146)} & Four-point population cross-term decay $\Gamma_4\to0$ over squarefree seeds, factored through CRT independence at four steps, for summable Borel--Cantelli densities toward DSL. & SCRTI gives cross-modulus independence at one time; FourPointPCV needs cross-time independence at one modulus. Four-time decay is HOD-type mixing (\#84), harder than CCSB; the four values are deterministic functions of one seed (\#115); pairwise PCV itself open. Superseded: false (\#156), DSL vacuous (\#160). \\
\midrule
\multicolumn{4}{@{}p{13.4cm}@{}}{\textsc{Aggregate gap} (AG, 7 entries) --- \emph{an average-case or aggregate bound does not give the per-fibre or per-class bound needed}.} \\[2pt]
101 & Bundle Walk / Weak MC / profinite {\footnotesize(S70)} {\footnotesize[revival 0]} & Prove Weak MC via the simultaneous walk on $\prod_{q\in M}(\mathbb Z/q)^\times$ using avoidance density $\delta(M)=\prod(q-2)/(q-1)\to0$; profinite reformulations generally. & BundleGap is weaker than SieveTransfer as a statement but no technique proves it otherwise; product-group characters do not factor (shared minFac index couples moduli), reducing to MMCSB; $\delta(M)\to0$ is population-level. In $(\mathbb Z/11)^\times$ cycling $\{3,4\}$ generates yet visits only $\{1,3\}$. \\
106 & VCB implies CCSB iff PED {\footnotesize(S78)} {\footnotesize[revival 2]} & Prove VCB directly via CRT+resampling, shift equivariance, ResamplingBound, fiber Parseval energy, and drop PED from VCB+PED$\Rightarrow$CCSB. & Telescope forces $(\mu-1)S_N=O(1)$: if $\mu\ne1$ CCSB is trivial, if $\mu\approx1$ kernel confinement kills CCSB, so (VCB$\Rightarrow$CCSB)$\iff$PED. Resampling error is \#98, deviation is \#90; fiber Parseval cross-character terms are unconstrained. Aggregate does not give per-fiber control. \\
121 & SMSB plus SE per-class escape {\footnotesize(S143)} {\footnotesize[revival 2]} & Use SecondMomentSquaredBound (global bad-set density) with SubgroupEscape to get per-class bad-set control (BadSetEquidistribution) and pointwise escape from class $-1$. & $|B|\le\delta N$ gives by pigeonhole some class with small bad density, not all and not $-1$; per-class uniformity needs $P(\mathrm{bad}\mid w=c)\approx P(\mathrm{bad})$, i.e. CME. SE gives generation, not statistics. Superseded: SMSB false at $\chi\equiv1$ (\#156). \newline {\footnotesize\color{leanink}$\checkmark$\,\code{marginal\_joint\_barrier\_witness}} \\
139 & Backward-dynamics chain broken everywhere {\footnotesize(S266)} {\footnotesize[revival 0]} & Chain ETA + DCTA + SRE $\Rightarrow$ CRTPropagationStep $\Rightarrow$ AEP $\Rightarrow$ DeathDensityLB $\Rightarrow$ SMLB $\Rightarrow$ LMG $\Rightarrow$ PositiveDensityRSD (Tao--Collatz backward density decomposition). & ETA false at $c=-1$ (\#136), DCTA false at $q=3$, SRE mis-stated (\#138), CRTPropagationStep false (absorption drains nonzero-class mass each step even with the corrected limit), AEP false (\#137), SMLB$(c)$ likely false since absorption forces $\mathrm{genSeq}$ to grow; chain zero-sorry but vacuous. \\
153 & Iwasawa / Euler-system receptacle {\footnotesize(S299)} {\footnotesize[revival 0]} & Iwasawa theory or Euler systems (Kolyvagin derivatives) with the EM orbit playing the role of the tower. & Kolyvagin derivatives need classes over the full squarefree lattice; the orbit supplies a single maximal flag $P_0 \mid P_1 \mid \cdots$. No $\mathbb{Z}_p$-tower (layer degrees are not $p$-powers), no motive, no period formula. \\
162 & Worst-case bag-prime bound in Lemma D {\footnotesize(S309)} {\footnotesize[revival 3]} & Bound the primes dividing the seed (the ``bag'' primes) in Lemma D by the worst case $\omega(m) \le \log_2 m$. & Once the threshold is type-determined, the worst-case bound is unbounded on a single type and destroys the uniform-in-$k$ constant: an average-case supply statement is being asked to hold per fibre. Repaired by a first moment, $\mathbb{E}\big[\sum_{r \mid m,\ r > z} 1/r\big] \le 2/(z\log z)$, and a Markov exclusion of a set of relative density $\le 1/C$ (\code{TailEstimate.seed\_divisor\_first\_moment}, \code{TailEstimate.markov\_divisor\_mass}). \\
173 & Applying \code{pathwise\_compensator} to the orbit of 2 {\footnotesize(S313)} {\footnotesize[revival 1]} & Instantiate the sure compensator at the true Euclid--Mullin orbit ($m=2$) to obtain an unconditional per-path statement about it. & Inapplicable at every horizon. The theorem requires $(\forall j<n,\ \tilde p_j \le Y)$ \emph{and} $\log Y \le n^2$ (and $e^{600}\le n$). No sure bound of that shape holds for the true orbit: the only sure size information permits $p_k \approx 2^{2^k}$, i.e. $\log Y \approx 2^n \gg n^2$. In the seed-average programme $Y$ is a \emph{policy} on a seed population and the violating seeds are absorbed into \code{ls\_plus}'s additive degenerate-tail term --- a single orbit cannot be absorbed into a tail term. Any future dispatch wanting to use the compensator on one orbit must first supply an unconditional $\log p_k \le k^2$, which is strictly stronger than (C$\infty$). \\
\midrule
\multicolumn{4}{@{}p{13.4cm}@{}}{\textsc{Collapse} (CO, 28 entries) --- \emph{the proposed hypothesis is definitionally (or by a short argument) equivalent to an existing one, usually CME, CCSB or MC itself}.} \\[2pt]
10 & WeakDecorrelation equals TailSE {\footnotesize(S11)} & Weaken DecorrelationHypothesis to ``not all multipliers in $\ker\chi$'' & Equivalent to TailSE, already known insufficient: gives escapes $R(N)\to\infty$ but not positive escape density, so no $o(N)$ walk sum \\
19 & BRE intermediate equals CCSB {\footnotesize(S14)} & Factor PED $\Rightarrow$ CCSB through BlockRotationEstimate (BRE) & BRE for $\chi$ of order $d$ is equivalent to CCSB for $\chi,\chi^2,\dots,\chi^{d-1}$ (time-equidistribution among $d$-th roots of unity), so BRE is CCSB, not a simplification \\
35 & BV plus threshold {\footnotesize(S17)} & Prove CCSB/DH for all but finitely many $q$ by averaging over moduli, finish with ThresholdHitting & All multi-modular variants collapse to single-modulus CCSB; BV exceptional set is density-zero, not finite; the EM-specific transfer is SieveTransfer \\
73 & SVE not provable from existing infrastructure {\footnotesize(S38)} & Prove SubquadraticVisitEnergy (excess energy $(p-1)\sum_a V_N(a)^2-N^2=o(N^2)$) from SE/TailSE, coprimality, PED, Dec+VdC or growth, then use \code{sve\_implies\_mmcsb}. & SE gives existence not density; coprimality invisible to the residue walk; PED constrains multipliers not walk phases; VdC at $h=1$ gives $\|S\|^2\le N^2/2$, a constant factor; growth is heuristic. Later proved SVE $\Leftrightarrow$ CCSB, so this collapses onto CCSB. \\
79 & HOD is not a useful attack target {\footnotesize(S40)} & Target HigherOrderDecorrelation as the hypothesis to prove, using the proved chain HOD + VdC $\to$ CCSB $\to$ MC. & HOD is strictly stronger than CCSB (\#76) with the same consequence, so proving HOD is the wrong direction; the reduction only documents the hierarchy and offers no simplification. \\
82 & Cyclic Littlewood--Offord / Tiep--Vu adaptation {\footnotesize(S46)} & Re-develop Tiep--Vu anti-concentration for $(\mathbb Z/q)^\times\cong\mathbb Z/(q-1)$ with weakened independence; state a MultiplierAntiConcentration Prop. & Reduces entirely to Dec/HOD: for a cyclic group the mechanism is $E[\prod\chi(X_k)]=\prod E[\chi(X_k)]$, independence used only to factor; MultiplierAntiConcentration $=$ PED; anti-concentration bounds the endpoint law, not the walk sum. No result exists for dependent products. \\
84 & Pseudo-independence notions reduce to Dec/HOD {\footnotesize(S46)} & Replace independence by pair mixing, exponential pair decay, $k$-point decay, block independence, or CME. & Pair mixing and exponential pair decay $=$ Dec (VdC $h=1$ gives $N/\sqrt2$); $k$-point decay $=$ HOD (unverifiable); block independence $=$ HOD at coarser scale; CME $=$ Dec only (\#81). EM feedback gives zero anti-concentration leverage. \\
85 & VdC plus EscapeDecorrelation equals QuadraticCCSB {\footnotesize(S49)} & For $d=2$ define EscapeDecorrelation ($h$-fold multiplier product sums $o(N)$) and prove VdC + EDH $\to$ QuadraticCCSB (Route C, formalized). & For $\chi^2=1$, $\chi(w(n+h))\chi(w(n))=\prod\chi(m(n+j))$, so EDH is literally ``walk autocorrelations are $o(N)$'', which is QuadraticCCSB by VdC; $h=1$ alone gives $O(N)$; BV supplies only $h=1$; PED does not imply EDH. \\
92 & Self-correction forcing equidistribution {\footnotesize(S59)} & Argue the growing excluded set (each new prime removed from future minFac candidates) feeds back and forces multiplier bias to zero, giving CME. & Feedback is summable ($\sum 1/\mathrm{seq}(n)<\infty$): it stabilises the walk but cannot force bias to zero. Necessary, not sufficient; extracting rates needs minFac conditioned on residue class, which is SieveTransfer/CME itself. Reformulation, not reduction. \\
93 & Density-1 CME / FEB equivalent to CME {\footnotesize(S61)} {\footnotesize[revival 0]} & Weaken CME to density-1 CME or an $L^2$ fiber-energy bound $CC_\chi(N)=\sum_a|F(a)|^2=o(N^2)$ as a strictly weaker intermediate between Dec and CCSB. & On the finite set of $q-1$ positions $L^2$ and $L^\infty$ coincide (Markov: $\#\{a:|F(a)|>\varepsilon N\}\le CC_\chi/(\varepsilon N)^2\to0$); removing a $\delta N$ fiber bias needs $\ge(\delta-\varepsilon)N$ deletions, exceeding the $o(N)$ budget. No $L^p$ interpolation; strictly stronger than CCSB. \newline {\footnotesize\color{leanink}$\checkmark$\,\code{cme\_implies\_feb}} \\
97 & LP avoidance framework equals DH {\footnotesize(S64)} & Linear program over Cesaro-limit measures on $(\mathbb Z/q)^\times$ with FPE/flow constraints; hope infeasibility for every avoidance set $T$ forces a hit of $-1$. & LP is feasible for many proper $T$ and always for $T=G\setminus\{-1\}$ (cannot detect single-element avoidance). Infeasibility for all $T$ with uniform FPE is equivalent to DynamicalHitting; only the FPE marginal is nontrivial and that is the open content. Fourier side reduces to CCSB. \\
99 & CME spectral gap / bias propagation {\footnotesize(S66)} & Direct CME attack via fiber flow-balance identity, CRT state-independence giving equal kernel rows, spectral gap decay, and return-product constraints $\prod\chi(m(j))=1$. & Flow-balance is one equation in $p-2$ unknowns (underdetermined); equal rows fail by biased subsequence selection; spectral gap is \#95 again; return products hold tautologically for any returning walk (group law), carrying no EM information. \\
100 & Walk periodicity dichotomy {\footnotesize(S69)} & Case-split on the walk being periodic (CCSB supposedly easy) versus aperiodic and attack each branch. & Periodic walks need not give $o(N)$ sums: period $p<q-1$ with nonzero one-period sum gives $\Theta(N)$. The trichotomy is a reformulation; ruling out periodicity removes the easy case and proving aperiodicity needs the same tools as CCSB. \\
103 & No third algebraic route to CCSB {\footnotesize(S72)} & Seek a decomposition of $S_N=\sum\chi(w(n))$ other than by value (fiber, CME) or lag (autocorrelation, HOD): Abel, differencing, Moebius, Dirichlet series, blocks, cross-modular CRT. & The telescope $\chi(w(n+1))=\chi(w(n))\chi(m(n))$ is the complete algebraic content. Abel gives $O(N^2)$ remainder; differencing yields norm-1 terms; Moebius needs multiplicativity in $n$; Dirichlet series needs structure. Blocks reduce to HOD (\#78), CRT to \#75/\#98/\#101. CME and HOD exhaust the routes. \\
104 & Summable Decorrelation collapses to VCB {\footnotesize(S74)} & Summable/Bottleneck Decorrelation (SD), stronger than VanishingConditionalBias, proposed to yield CCSB together with PED. & VCB+PED$\Rightarrow$CCSB was proved outright (\code{vcbPedImpliesCcsb}), so SD is strictly stronger than VCB yet adds zero leverage; verifying SD for EM meets \#90/\#98. Not a published concept. \\
105 & First passage / ExistentialCME equals DH {\footnotesize(S75)} & Exploit that MC needs one hit of $-1$: phased Fourier identity $V_N(-1)(q-1)=N+\sum_{\chi\ne1}\chi(-1)^{-1}S_N(\chi)$ and ExistentialCME. & The phased identity is character orthogonality restating $V_N(-1)=0$; ExistentialCME ($\exists c,n$: $w(n)=c$, $m(n)=-c^{-1}$) is literally DH; inter-character cancellation needs HOD (\#79). A walk can generate the group, use every step cofinally, and avoid one element forever. \\
110 & Transition matrix / Doeblin equals CME {\footnotesize(S81)} {\footnotesize[revival 0]} & Recast CME as convergence of the empirical transition matrix $K_N(a,b)/V(a)$ to uniform and use Markov tools (spectral gap, mixing, Doeblin, non-homogeneous chains). & $K_N(a,b)/V(a)\to$ uniform is CME (\code{cme\_iff\_transition\_char\_vanish}); DoeblinConvergenceForEM $=$ CME by \code{rfl}. Transition counts are joint counts, so the convolution-kernel identity itself needs CME. All convergence techniques need randomness or stationarity. \newline {\footnotesize\color{leanink}$\checkmark$\,\code{doeblin\_eq\_cme}} \\
112 & Order-3 Moebius death function {\footnotesize(S83)} & Exploit the order-3 Moebius map $b(a)=(1-a)^{-1}$ on forbidden multipliers ($m_1m_2m_3=-1$) as leverage against DH. & It constrains the death-curve geometry, not the walk dynamics; the dynamically relevant map is the order-2 involution $f(c)=-c^{-1}$, already formalized. Avoidance graph is a perfect matching with no combinatorial leverage; Marginal/Joint Barrier untouched. \\
114 & Missing-prime accumulation / Borel--Cantelli {\footnotesize(S97)} & Second-moment / Borel--Cantelli / Kochen--Stone argument that a missing prime $q$ accumulates death-channel opportunities along the orbit. & Pairwise death-channel independence is CME restricted to one fiber; the self-consistent avoidance model restates the earlier orbit analysis; Kochen--Stone needs pairwise quasi-independence of $\{q\mid P(n)+1\}$, the joint-distribution barrier. Population-to-orbit transfer is SieveTransfer (\#90/\#98/\#107). \\
120 & SelfCorrectingDrift equivalent to SVE {\footnotesize(S142)} {\footnotesize[revival 2]} & Lyapunov route SCD ($R(N)=o(N^2)$) $\Rightarrow$ VE $\Rightarrow$ SVE $\Rightarrow$ MC, hoping mixing-time or concentration arguments bypass CME. & \code{lyapunov\_telescope}: $L(N)=2R(N)+N(q-2)/(q-1)$, so $R=o(N^2)\iff L=o(N^2)$ and SCD is SVE restated in drift language (above threshold $Q_0$); maps to \#73/\#92/\#95. No new leverage. \newline {\footnotesize\color{leanink}$\checkmark$\,\code{lyapunov\_telescope}} \\
122 & Tail Window Decorrelation {\footnotesize(S144)} & TWD: split multiplier character sums into non-overlapping windows of length $K$, require cross-window terms to vanish, feed into CCSB. & TWD controls multiplier sums, giving Dec at best, and Dec does not imply CCSB (\#20/\#58/\#117); the walk-sum version is block HOD at scale $K$ (\#84) $=$ CME. Cycling $\{g,g^{-1}\}$ has bounded cross terms but walk sum $\Theta(N)$. Superseded: TWD false (\#156). \\
124 & T5.7 Lyapunov--fiber coupling {\footnotesize(S154)} & Couple visit deviations with fiber sums, $J(N)=\sum_a\mathrm{visitDeviation}(a)\,\mathrm{fiberMultCharSum}(a)$, and use its recurrence or Cauchy--Schwarz for DSL/CME leverage. & $J$ is marginal (bilinear contraction loses joint information); the one-step recurrence contains the active-fiber sum $F(w(N),N)$, the CME gap; Cauchy--Schwarz bounds go the wrong way; $J=o(N)$ is one scalar constraint on a $(q-1)$-dimensional problem. Maps to \#110/\#104/\#120. \\
143 & Congruence certificate collapses to MC {\footnotesize(S298)} {\footnotesize[revival 0]} & Hypothesis \code{IC\_min}: a missing prime is certified by a propagating congruence invariant, in analogy with the max sequence (CvdP/Booker). & The No-Invariant Theorem \code{no\_cvdp\_obstruction} (eviction, free-state fullness via Dirichlet, CRT-reach) shows no propagating congruence set blocks a missing prime, so \code{IC\_min} $\leftrightarrow$ \code{MullinConjecture}; \code{SingleHitHypothesis}/\code{DynamicalHitting} make it vacuous. Any disproof of MC must be non-congruential. \newline {\footnotesize\color{leanink}$\checkmark$\,\code{CvdP.ic\_min\_network}} \\
145 & Simultaneous avoidance decouples {\footnotesize(S298)} {\footnotesize[revival 1]} & Heath-Brown-genre pairwise incompatibility: two or $k$ primes cannot all be missing, using the shared finite prefix that Finite Hitting forces on all hits. & $k$-fold missingness is the monotone condition $S\cap\mathrm{Im}(\minFac)=\emptyset$ with no interaction term; the apparent tail confinement is unconditional (holds for present primes too); the only shared budget (rogue characters + large sieve) is mispriced by $2^N/N$ (\#108); embeddings give almost-all statements (\#90). \newline {\footnotesize\color{leanink}$\checkmark$\,\code{CvdP.hittingSet\_finite}} \\
149 & Ledger small-prime output weaker than injectivity {\footnotesize(S299)} {\footnotesize[revival 0]} & Bound total small-prime consumption $\sum_{p\le y} h_\infty(p)$ via the hitting ledger (\code{hitting\_ledger\_bound}). & The ledger gives $\sum_{p\le y} h_\infty(p) \le \pi(y)^2$, strictly weaker than the injectivity bound $\pi(y)$ (\code{seq\_injective}), because ``$\exists p\le y$ with $p\mid N_n$'' is ``$\minFac N_n \le y$'': the roughness trap, quantified. \newline {\footnotesize\color{leanink}$\checkmark$\,\code{seq\_injective}} \\
150 & Covering-system congruence obstructions {\footnotesize(S299)} {\footnotesize[revival 0]} & Escape the single-modulus No-Invariant Theorem with an Erdős-style covering family of moduli. & Covering systems are finite by definition, so \code{exists\_tail\_coprime} at $m=\prod T$ kills fixed-finite-prime-set covering (\code{no\_finite\_prime\_covering}); assembling at the lcm gives one set at one modulus and \code{no\_cvdp\_obstruction} is set-generic. Residual escapes: unbounded/profinite families and anatomy invariants. \newline {\footnotesize\color{leanink}$\checkmark$\,\code{CvdP.no\_covering\_family\_obstruction}} \\
159 & $\varepsilon$-interpolation family reformulates MC {\footnotesize(S307)} {\footnotesize[revival 0]} & The $(1-\varepsilon)\minFac + \varepsilon\cdot$random family: prove almost-sure capture at every fixed $\varepsilon>0$ and pass $\varepsilon\to 0$ to get MC. & \code{mullin\_iff\_exists\_failWeight\_bound}: for $q\neq 2$, MC at $q$ iff some $(\varepsilon,N,c)$ with $N\varepsilon<c$ has $\mathrm{failWeight}\le 1-c$; outside $N<c/\varepsilon$ no bound exists unless MC holds; if $q$ never occurs, $\mathrm{failWeight}\ge(1-\varepsilon)^N$, compatible with a.s. capture at every fixed $\varepsilon$. MC relocated into a finite window, not weakened. \newline {\footnotesize\color{leanink}$\checkmark$\,\code{mullin\_iff\_exists\_failWeight\_bound}} \\
169 & Harmonic charge budget applied to a single orbit {\footnotesize(S313)} {\footnotesize[revival 1]} & Use \code{LargeStepRoughness.charge\_sum\_le\_harmonic} --- a sure, per-path, hypothesis-free inequality --- to constrain the primes missed by one Euclid--Mullin orbit (in particular the true orbit, seed $m=2$). & Vacuous: the budget is an \emph{identity}, not an inequality with slack. The box starts at $r-1$, a charge decrements it by \emph{exactly} one (\code{boxCard\_succ\_of\_charged}) and a non-charge leaves it unchanged (\code{boxCard\_of\_not\_charged}), so the charged box sizes are exactly $r-1,\dots,r-C$ and $\sum_{\text{charged}} 1/|B| = H_{r-1}-H_{r-1-C}$. The theorem is therefore \emph{equivalent}, for a single path, to $C \le r-2$ --- and the Lean proof is transparently that, an injection into \code{Finset.range (r-1)}. But $C$ counts distinct cofactor residues at $r$-exposed steps, which are units mod $r$, so $C \le r-2$ is forced by counting alone. Zero orbit-distinguishing content at every $r$ and every horizon; the aggregate form is within an absolute constant of generic behaviour. The reading ``$\approx r\log r$ declines are permitted'' is a lossy relaxation --- only \emph{charges} are bounded, declines are not. \\
\midrule
\multicolumn{4}{@{}p{13.4cm}@{}}{\textsc{Circularity} (CI, 13 entries) --- \emph{the argument presupposes the conclusion or an equivalent}.} \\[2pt]
12 & PairDecorrelation is circular {\footnotesize(S11)} & Control all lag-$h$ shift correlations of the walk (PairDecorrelation) and apply van der Corput & $h=1$ correlation is just the multiplier sum (\code{walk\_shift\_one\_correlation}); $h \ge 2$ requires PairDecorrelation, which is equivalent to walk equidistribution / DH itself — circular \\
14 & Expander graphs for PED {\footnotesize(S11)} & Expander mixing on a Cayley-type graph of the multipliers & Circular: PED is needed as input to get expansion; $(\mathbb{Z}/q)^\times$ abelian, quasirandomness 1, Bourgain--Gamburd inapplicable \\
16 & Selberg sieve on EM products {\footnotesize(S13)} & Apply Mathlib \code{SelbergSieve}/\code{BoundingSieve} to bound the minFac distribution of $\mathrm{prod}(n)+1$ & Density approximation $\mathrm{multSum}(d) \approx \nu(d)\cdot\mathrm{mass} + \mathrm{rem}(d)$ needs $\#\{n : d \mid \mathrm{prod}(n)+1\}$, i.e. walk equidistribution mod $d$ — the very thing to be proved. Population use later itself false (\#160) \\
21 & Weyl-criterion shortcut for partial products {\footnotesize(S15)} & Use the Weyl criterion on $\mathbb{Z}/d$ for the walk $P_n = \omega^{s_n}$ to bypass CCSB & Needs $s_n$ equidistributed mod $d$; van der Corput / ergodic tools reduce to PairDecorrelation or stronger; the finite Weyl criterion (later proved) is CCSB itself \\
52 & Five direct CCSB approaches via product structure {\footnotesize(S22)} & Q1 integer divisibility $m\mid P(n)+1$ at character level; Q2 profinite walk with minimality coupling; Q3 Hooley sieve for $\minFac$; Q4 entropy/variance; Q5 prime density. & Q1: divisibility invisible to characters. Q2: reduces to SieveTransfer/CME. Q3, Q5: reach only PED, insufficient for $d\ge3$. Q4: walk equidistribution is CCSB by orthogonality (circular). All single-modulus routes reach PED at best. \\
54 & BV decomposition obstacles, CrossPrime, direct LoD {\footnotesize(S25)} & Decompose \code{BVImpliesMMCSB} into sequence-vs-population, $L^2$-vs-pointwise, counting-vs-walk obstacles; propose CrossPrimeAmplification via minimality and a direct EM level of distribution. & BVImpliesMMCSB faces the same generic-to-EM gap as SieveTransfer; CrossPrimeAmplification has no published technique; a direct LoD proof is equivalent to SieveTransfer, and \code{EMHasLevelOfDistribution} as stated has an exponentially large bound (later \#96). \\
61 & ArithLS to MMCSB via occupation measure {\footnotesize(S35)} & Apply the arithmetic large sieve (or \code{gct\_full\_char\_sum\_bound}) with coefficients $a_n=\#\{k:\ w(k)=n\}$ to derive MultiModularCSB. & Circular: $\sum|a_n|^2$ is the $L^2$ norm of the visit distribution, large exactly when the walk is not equidistributed; dense coefficients give the trivial bound $N^2(p+1)$, worse than the triangle inequality. \\
113 & Cycle product equidistribution circular {\footnotesize(S91)} & Equidistribute products $R_k$ of multipliers over walk cycles/return times mod $p$, hoping product structure ($\ell\ge2$) helps. & By the telescope $\chi(R_k)=a_{k+1}/a_k$, equidistribution of $R_k$ is a walk character-sum bound at return times, i.e. CCSB for $p$; lag-1 autocorrelation is \#92/\#98, cross-cycle decorrelation \#98/\#58, $\ell\ge2$ reduces to HOD (\#79). \\
116 & Sieve-theoretic transfer for DSL circular {\footnotesize(S114)} {\footnotesize[revival 0]} & DSL sub-approach D: prove EMDirichlet mod $q$ with a Selberg-type sieve over auxiliary primes $r$. & The sieve axiom $\omega(r)\sim1/r$ is EMDirichlet mod $r$, so the argument assumes EMDirichlet for all auxiliaries: circular for $q\le L$, reduces to BVImpliesMMCSB for $q>L$. CRT independence is a scope error (for fixed $r$, $r\equiv a$ mod $q$ is deterministic). Superseded: DSL vacuous (\#160). \\
132 & L-function factorization is circular {\footnotesize(S173)} {\footnotesize[revival 1]} & Factor $L(s,\chi) = L_{\mathrm{EM}}(s,\chi)\,L_{\mathrm{non\text{-}EM}}(s,\chi)$ and use zero-freeness of the EM factor on $\Re s = 1$ for equidistribution of EM primes. & Controlling $L_{\mathrm{non\text{-}EM}}$ requires knowing which primes are non-EM, i.e. MC itself; the reformulation reduces to the hypothesis it aims to prove (maps to \#90). \\
134 & No Tauberian lever for $L_{\mathrm{EM}}$ {\footnotesize(S173)} {\footnotesize[revival 1]} & Apply Tauberian theorems to the EM-specific Dirichlet series $\sum (\prod n)^{-s}$ / $L_{\mathrm{EM}}(s,\chi)$. & $\prod n \ge 2^{n+1}$ makes the series converge for all $s>0$ (\code{accum\_reciprocal\_summable}); $L_{\mathrm{EM}}$ is entire on $\Re s>0$ with no pole at $s=1$. The standard PNT route uses $L(s,\chi)$ over all primes, already formalized. \\
166 & ``The bag has caught up'' justification of $S_k \asymp 1$ {\footnotesize(S309)} {\footnotesize[revival 3]} & Justify the linear growth of the compensator by the heuristic that only primes $r \gtrsim k/\log k$ are still active at depth $k$, the smaller ones having already been captured (\S1.7). & Circular: it assumes the small primes have been captured, which is precisely the population form of the conclusion (LS)/Theorem~C. Any proof of (LS) routing through ``the bag has caught up'' presupposes its own conclusion. The non-circular route, and the one actually formalized, is the harmonic charge budget $\sum_{r<2n} H_{r-1} \le \theta(2n)+\pi(2n) = O(n)$, which uses only that an \emph{uncaptured} prime pays $1/|B|$ into its own finite budget whenever it declines (\code{LargeStepRoughness.charge\_sum\_le\_harmonic}). \\
170 & ``The box has not collapsed'' as a usable branch {\footnotesize(S313)} {\footnotesize[revival 0]} & Extract information about a missed prime $r$ from the surviving branch of the missed-prime dichotomy, namely that the box $B_k(r)$ stays large / never empties. & Circular. \code{seed\_mem\_box}'s Lean proof is literally ``exposure $\Rightarrow$ non-divisibility'' (it ends in \code{not\_dvd\_succ\_of\_exposed\_avoid}), so $|B_k(r)|\ge 1$ unwinds to ``$r \nmid \Prod{m}(j) + 1$ at every $r$-exposed $j<k$'', i.e. the walk mod $r$ avoids the death class --- which is $\neg\mathrm{DynamicalHitting}(r)$, the hypothesis one is trying to contradict. The other branch (``$r$ rarely exposed'') is refuted outright by \code{few\_small\_multipliers}: at least $n-\pi(r)$ of the first $n$ steps are $r$-exposed regardless. Nor can the circle be broken by weakening: the only nonvacuous quantitative negation, $C \ge r-1$, \emph{is} MC($r$). Exact mirror of \#166 --- \#166 blocks the entrance to (LS), this blocks the exit. \\
\midrule
\multicolumn{4}{@{}p{13.4cm}@{}}{\textsc{Structurally false} (SF, 30 entries) --- \emph{the proposed statement is false: an explicit counterexample or a proved refutation}.} \\[2pt]
6 & Direct SE implies DH (generation not coverage) {\footnotesize(S5)} & Prove DynamicalHitting directly from SubgroupEscape / generation of $(\mathbb{Z}/q)^\times$ by the multipliers & Counterexample: in $\mathbb{Z}/6$ steps $2,2,2,\dots$ generate but partial sums cycle $0,2,4$ and miss all odd elements; generation does not imply coverage. Later $\mathbb{Z}/4$ witness \\
20 & Decorrelation does not imply CCSB {\footnotesize(S15)} {\footnotesize[revival 1]} & Prove ComplexCharSumBound directly from multiplier decorrelation (step character sums $o(N)$) & Counterexample: increments equidistributed in $\mathbb{Z}/3$ (bunched by type) but walk stalls; alternating steps in $\mathbb{Z}/4$ generate yet the walk misses 2. Step cancellation is not walk cancellation \newline {\footnotesize\color{leanink}$\checkmark$\,\code{alternating\_walk\_misses\_two}} \\
36 & PED to BRE fails for $d\ge3$ {\footnotesize(S18)} & Prove BlockRotationEstimate (PED $\Rightarrow$ CCSB) from positive escape density & Counterexample on $\mathbb{Z}/3$: block values $1,\omega,1,\omega,\dots$ with lengths $M,1,M,1,\dots$; escape density $2/(M+1) > 0$ but walk sum $\approx N$. Escape density fixes frequency, not distribution among $d-1$ rotations \\
37 & NoLongRuns plus PED implies BRE fails {\footnotesize(S18)} & Add NoLongRuns($L$) to PED to control block lengths & Variable block lengths still align adversarially with character phases for $d \ge 3$; Abel gives total variation $\le 2N$, bound $O(N)$ = trivial. Order-2 case provable (\code{bre\_order2\_from\_noLongRuns}) \\
42 & Cumulative SieveEquidist implies NoLongRuns {\footnotesize(S20)} & Prove \code{SieveEquidistImpliesNoLongRuns} from cumulative-density \code{SieveEquidistribution} & Cumulative density on $[0,N)$ does not exclude a late run of $L$ kernel steps at $n \gg L$; fixed by window version \code{StrongSieveEquidist} (\code{strongSieveEquidist\_noLongRunsAt} proved). \#160 later voided the sieve-route endpoint \\
45 & DPED to CCSB fails for $d\ge3$ {\footnotesize(S21)} & Distributional PED (every non-identity $\chi$-value with positive frequency) as PED $\Rightarrow$ CCSB intermediate & Counterexample: rotations alternate $\omega,\omega^2$ (DPED holds), cumulative products cycle $\omega,1$, walk sum $(N/2)(1+\omega)$ of norm $N/2 = \Theta(N)$; DPED controls pointwise not sequential correlations; PED $<$ DPED $<$ RI $\approx$ CCSB \\
46 & Escape-value distribution intermediates below CCSB {\footnotesize(S21)} & Any hypothesis on the pointwise distribution of escape values $\chi(m(n))$ (PED, DPED, balanced escapes) as a bridge PED $\to$ CCSB for order $d\ge 3$. & Sequential correlations matter, not pointwise statistics: rotations alternating $\omega,\omega^2$ have perfect pointwise balance yet cumulative products cycle $1,\omega$, giving walk sum $(N/2)(1+\omega)=\Theta(N)$. Only controlling cumulative products works, which is CCSB itself. \\
76 & CCSB does not imply HOD {\footnotesize(S40)} & Hoped-for equivalence between walk-sum cancellation (CCSB) and $h$-fold multiplier autocorrelation decay (HigherOrderDecorrelation). & CCSB bounds $\sum\chi(w(n))$, HOD bounds $\sum\prod_{j<h}\chi(m(n+j))$; an equidistributed walk with anti-correlated consecutive increments satisfies CCSB and violates HOD at $h=2$. Hierarchy Dec $<$ HOD $\to$ CCSB $\Leftrightarrow$ SVE. \\
83 & Inverse Littlewood--Offord for products false {\footnotesize(S46)} & ``Large walk sum implies multiplier concentration in $\ker\chi$'' (Tao--Vu torsion-implies-dispersion intuition) to turn CCSB failure into structure. & False for $d\ge3$: $\chi$ of order 3, multipliers alternating $\omega,\omega^2$: walk visits $\{1,\omega\}$, $|S_N|=N/2$ yet no multiplier lies in $\ker\chi$. The inverse direction is just the contrapositive of CCSB. \\
87 & General escape-sum bound fails, $d\ge3$ {\footnotesize(S54)} & Extend \code{escape\_sum\_order2\_bounded} (unweighted escape sum $O(1)$ for order-2 $\chi$) to order $d\ge3$ to generalize the kernel-block reduction CCSB $\leftrightarrow$ kernel-CCSB. & For $d=2$ every escape rotation is $-2$, so the telescoped weighted sum factors; for $d\ge3$ rotations $\chi(m)-1$ take $d-1$ values and $\|\sum z_nr_n\|\le2$ does not bound $\|\sum z_n\|$: adversarial phase alignment gives $\Theta(N)$. Confirms \#36. \\
88 & d=2 block-length balance from PED alone {\footnotesize(S55)} & Use the $d=2$ alternating series of kernel-block lengths (every escape negates the walk character) plus PED to close QuadraticCCSB. & PED controls total kernel density, not the distribution of kernel steps across blocks; alternating long/short blocks $L_k=\alpha N/K$, $(2-\alpha)N/K$ give alternating sum $\Theta(N)$ with $\Omega(N)$ blocks. Balance $\Leftrightarrow$ translated QR/QNR symmetry $\Leftrightarrow$ localized SieveTransfer. \\
107 & Bottleneck Decorrelation axioms insufficient {\footnotesize(S79)} & Abstract theory: Per-Step CRT Bottleneck, Exponential Growth and Generation axioms for a class of walks should imply VCB. & Counterexample on $(\mathbb Z/3)^\times$: block-structured multipliers ($N=4K$, escape blocks $m=2$, kernel blocks $m=1$) satisfy all axioms but $F(1)/V(1)=1/3$, $F(2)/V(2)=-1$, no common $\mu$. Per-step CRT is pointwise; VCB needs aggregate joint (position, multiplier) control. \\
119 & Substitution Principle false {\footnotesize(S140)} & SP: pairwise coprime $N_i\equiv c$ mod $q$ with distinct minFac$(N_i)$ should have $\sum\chi(\mathrm{minFac}\,N_i)=o(I)$; apply at fiber return visits. & Counterexample $q=5$: $N_i=p_ir_i$ with distinct primes $p_i\equiv2$, $r_i\equiv3$ mod 5, $r_i>p_i>2^{2^i}$; minFac distinct, $N_i\equiv1$, coprime, yet the sum is $I\chi(2)=\Theta(I)$; generalizes to all $q\ge3$. SP-for-EM equals CME. \newline {\footnotesize\color{leanink}$\checkmark$\,\code{sp\_eq\_cme}} \\
125 & Pairwise is not $k$-wise (XOR) {\footnotesize(S164)} {\footnotesize[revival 3]} & Derive UniformProfiniteEquidist ($k$-wise modular decorrelation) from CCSB plus pairwise cross-modulus decorrelation CPD for the profinite Weyl criterion. & XOR counterexample: unit-modulus $X_1,X_2$ with pairwise cancellation and $X_3=X_1X_2$ give $\sum X_1X_2X_3=N$. Cauchy--Schwarz/Hoelder are trivial, VdC reduces to multiplier CPD, Tao--Teravainen needs multiplicativity; UPE needs $k$-wise for all $k$ as a primitive. \\
129 & FFLM false: cyclotomic counterexample {\footnotesize(S168)} {\footnotesize[revival 3]} & FF-EM monodromy chain: large Galois groups of $\mathrm{ffProd}(n)+1$ (FFLM) plus Deligne equidistribution to obtain orbit-specific FF-CME. & Over $\mathbb{F}_2$, $\mathrm{ffProd}(2)+1 = \Phi_5(t)$ with $\mathrm{Gal} = \mathbb{Z}/4$ (abelian; FF-EM products divide $t^{p^d}-t$, so Galois groups are abelian). Deligne is a family statement (sequential gap = \#90), and Frobenius cycle type does not determine $\minFac \bmod Q$. \\
130 & Generation does not imply coverage {\footnotesize(S170)} {\footnotesize[revival 1]} & Selection-bias-neutral / conditional character equidistribution to turn population multiplier CCSB into FF-multiplier CCSB; more generally subgroup escape as a route to hitting. & \code{ConditionalCharEquidist} $=$ FF-CME $=$ CME definitionally (collapse to \#90); and generation is not coverage: steps $\{1,3\}$ generate $\mathbb{Z}/4$ but the alternating walk visits only $\{0,1\}$, so subgroup escape never forces the walk to hit $-1$. \newline {\footnotesize\color{leanink}$\checkmark$\,\code{alternating\_walk\_misses\_two}} \\
136 & ETA false for death class $c=-1$ {\footnotesize(S265)} {\footnotesize[revival 2]} & \code{EnsembleTransitionApprox} (Tao-style backward dynamics): ensemble transition probability from accumulator class $c$ to multiplier class $b$ tends to $1/(q-1)$. & When $\mathrm{genProd} \equiv -1 \pmod q$, a positive fraction ($\sim C_1/\log q$, Mertens) have $\minFac(\mathrm{genProd}+1) = q$, i.e. $b=0$; mass leaks so $T(-1,b) < 1/(q-1)$ for $b\neq 0$. At $q=3$ the split-off DCTA is provably false (\code{genSeq\_eq\_three\_of\_genProd\_mod3}). \\
137 & AEP false at $q=3$ (absorption) {\footnotesize(S266)} {\footnotesize[revival 1]} & \code{AccumulatorEquidistPropagation}: nonzero classes of $\mathrm{genProd}\,n\,k \bmod q$ keep density $1/(q-1)$ (later $r/(r^2-1)$) as $k$ grows. & At $q=3$ death ($\mathrm{genProd} \equiv 2$) forces $\mathrm{genSeq}=3$, hence $\mathrm{genProd}\equiv 0 \pmod 3$ forever; $F_k(1),F_k(2) \sim C 2^{-k} \to 0$ while $F_k(0)\to 1$. Heuristically false at every fixed $q$ (absorption rate $\sim 1/\log q$ per step); \code{DeathDensityLB}$(3,c)$ false for all $c>0$. \newline {\footnotesize\color{leanink}$\checkmark$\,\code{death\_then\_permanent\_absorption}} \\
138 & SRE limit wrong: $r/(r^2-1)$ {\footnotesize(S266)} {\footnotesize[revival 1]} & \code{SquarefreeResidueEquidist}: density of squarefree $n \equiv a \pmod r$ among squarefree $n$ tends to $1/(r-1)$, as base case of the CRT-propagation chain. & For $\gcd(a,r)=1$ the density is $(1/r)\prod_{p\neq r}(1-p^{-2})/\prod_p(1-p^{-2}) = r/(r^2-1)$; class $0$ has density $1/(r+1) \neq 0$; discrepancy factor $r/(r+1)$ (for $r=3$: $3/8$ vs $1/2$). The stated Prop was a false open hypothesis; limit corrected in code. \\
140 & Cauchy--Davenport coverage vacuous ($\mathrm{minOrder}=2$) {\footnotesize(S234)} {\footnotesize[revival 1]} & Coverage of $(\mathbb{Z}/q)^\times$ by iterated products of factor sets via Mathlib's \code{cauchy\_davenport\_minOrder\_mul}, extended by Kneser. & For every prime $q\ge 3$, $-1$ has order $2$, so $\mathrm{minOrder}(\mathbb{Z}/q)^\times = 2$ and CD gives only $|AB|\ge 2$; the safe-prime hypothesis forces $q=3$; Kneser (not in Mathlib) would still face ordering (\#4) and generation$\neq$coverage (\#130). \newline {\footnotesize\color{leanink}$\checkmark$\,\code{iterated\_product\_dead\_end\_landscape}} \\
142 & Mod-4 CvdP argument closes only $q=5$ {\footnotesize(S298)} {\footnotesize[revival 0]} & Generalize the Cox--van der Poorten mod-4 omission argument for the max sequence to further primes and transport a propagating congruence obstruction to the min sequence. & For $q=11,13$ the forced smooth sets $5^a7^b11^c(13^d)$ meet the class $3 \bmod 4$ since $7\equiv 11\equiv 3$; on the min side $\minFac N=q$ is an open congruence condition: $N=35$ satisfies all max-side congruence hypotheses, has $\minFac 35=5$, is not a $5$-power. \newline {\footnotesize\color{leanink}$\checkmark$\,\code{min\_side\_no\_smoothness\_forcing}} \\
146 & Zero-configuration (eventually-prime) branch {\footnotesize(S299)} {\footnotesize[revival 2]} & Consumption/ledger receptacles for MC; two agents independently hit the branch where $P_n+1$ is prime for all large $n$. & Under perpetual primality $P_{n+1}=P_n(P_n+1)$ and the walk mod $q$ is autonomous, $w\mapsto w^2+w$; $w^2+w+1$ has no root in $\mathbb{F}_q$ for $q\equiv 2\pmod 3$, so MC would fail on a density-$1/2$ set. Consumption receptacles are vacuous there ($\omega(P_n+1)-1\equiv 0$); ledgers must be gated on (C$\infty$). \newline {\footnotesize\color{leanink}$\checkmark$\,\code{AutonomousBranch.eventually\_prime\_implies\_not\_mullin}} \\
151 & Confinement cohomology Gap is false {\footnotesize(S299)} {\footnotesize[revival 0]} & Treat the Fact-4 confinement box (avoidance subcomplex) cohomologically and show $H^0$ vanishes. & \code{free\_transition} with \code{exists\_tail\_coprime} makes every tail state free, so the box is forward-mobile under the over-approximated \code{Transition}; $H^0\neq 0$. Restricting to orbit-realizable edges is DSL. Side-finding: \code{free\_transition} is rule-symmetric; the min/max break point is capture (\code{forcingState\_captures}), not fullness. \newline {\footnotesize\color{leanink}$\checkmark$\,\code{CvdP.free\_transition}} \\
156 & Uncentered ensemble character layer false {\footnotesize(S307)} {\footnotesize[revival 0]} & Ensemble variance route: \code{StepDecorrelation}, \code{CharSumVarianceBound}, \code{EnsembleCharSumConcentration}, \code{FourPointPCV}, \code{SecondMomentSquaredBound}, TWD, quantified over all $\chi$ with only $\mathrm{normSq}\le 1$. & Take $\chi\equiv 1$: \code{ensembleAvg} of $1$ is $1$, so the SD limit is $1$ not $0$; energy $K^2 \not\le CK$; bad density $1$; four-point average $\equiv 1$; $E[E^2]=K^4$. The side conditions ($\chi(0)=0$, $\sum\chi=0$) of \code{MultCancelToWalkCancel} were never back-ported. Repair: centered per-$\chi$ covariance plus Cesàro drift. \newline {\footnotesize\color{leanink}$\checkmark$\,\code{UncenteredRefutations.not\_stepDecorrelation}} \\
157 & Fixed-step multiplier equidistribution false {\footnotesize(S307)} {\footnotesize[revival 0]} & \code{EnsembleMultiplierEquidist} (limit $1/(q-1)$) and \code{JointStepEquidist} for $\mathrm{genSeq}\,n\,k \bmod q$ over squarefree starting points $n\le X$. & $\mathrm{genSeq}\,n\,0=2$ for every odd $n$ (\code{genSeq\_zero\_of\_odd}) and at least half the squarefree $n$ are odd, so mass $\ge 1/2$ sits on $2\bmod q$: false for $q\ge 5$. Not a parity artifact: on the family $2p$ the Dirichlet density of $\{\minFac(2p+1)=3\}$ is $1/2$. Independent of absorption. \newline {\footnotesize\color{leanink}$\checkmark$\,\code{UncenteredRefutations.not\_ensembleMultiplierEquidist}} \\
160 & PE/MFRE/RoughLPFEquidist false: head domination {\footnotesize[revival 0]} & Population layer ANT $\Rightarrow$ MFRE $\Rightarrow$ PE $\Rightarrow_{\mathrm{DSL}}$ CME: the class $\minFac\equiv a\pmod q$ among $q$-rough integers has density $1/(q-1)$ ``by Dirichlet''. & The density of $\{\minFac m=p\}$ is $w_p=p^{-1}\prod_{r<p}(1-1/r)$; weights telescope, so the class density is the convergent series $\sum_{p\equiv a}w_p$ (\code{tendsto\_classCount\_div}) and RoughLPFEquidist is the identity $\sum_{p\equiv a}w_p=c_q/(q-1)$, on which Dirichlet is silent: the class of the least prime above $q$ receives more than its share (about twice for large $q$; the excess at any fixed $q$ is a finite positive-term computation, deliberately not run in Lean). Same for MFRE/PE, UCE. \newline {\footnotesize\color{leanink}$\checkmark$\,\code{HeadDomination.roughLPFEquidist\_iff}} \\
164 & Block-chaining substitute for Freedman {\footnotesize(S309)} {\footnotesize[revival 3]} & Replace a finite-horizon Freedman inequality by chaining a martingale estimate over blocks of consecutive steps (\S F.5.5 of the (LS) proof sketch). & False on two independent counts: (i) the harmonic charge budget is \emph{global} and admits no per-block version --- a middle block may consist entirely of high-exponent steps, so the per-block lower bound fails; (ii) the goodness of step $k$ depends on $\tilde p_{k+1}$, so the block indicator is not block-past-measurable and the chaining is not a martingale. Replaced by a finite-tree exponential supermartingale driven by the \emph{sure} compensator bound (\code{TreeChernoff.chernoff\_quarter\_local}), which needs neither stopping nor blocks. \\
165 & Omitting $r \ne q$ from the $q$-free box process {\footnotesize(S309)} {\footnotesize[revival 1]} & Run the box/charge process of the (LS) proof over all primes $r \le Y$, including $r = q$ itself. & False at $r = q$: the $q$-free selector $\tilde p = \minFac_{\ne q}$ ignores $q$ entirely, so the brink lemma (F3) fails at $r = q$ and the step survival $S_k$ computes the wrong probability. The $q$-free type is a function of $m \bmod M_Y$ with $q \nmid M_Y$; $q$ is not a box coordinate at all, and the $q$-coordinate is carried separately by the capture identity. \\
171 & Sure layer can force a capture ($q$-free model obstruction) {\footnotesize(S313)} {\footnotesize[revival 0]} & Derive, from the sure layer alone (charge budget, brink lemma, distinctness, \code{chargeBudget\_le}, \code{pathwise\_compensator}), that some prescribed prime $q$ must be captured by the orbit. & Structurally impossible. Every theorem of the sure layer is proved \emph{about} $\code{genSeqAvoid}\,q\,m$, and \code{SeedCapture.genSeqAvoid\_ne\_avoided} proves that this dynamics \emph{never selects} $q$. So at every finite horizon on which the $q$-free orbit is nondegenerate, the entire layer is satisfied by a dynamics that misses $q$ by construction; no consequence of the layer alone can force any capture. Stated precisely so as not to overclaim: capture \emph{is} recoverable per-orbit via \code{captured\_iff\_mem\_visited} at $m'=m$, but nothing sure forces \code{visitedSetAvoid} to grow so as to contain a \emph{prescribed} class. The single cleanest reason the orbit direction is closed, and the only one with a ready-made formal witness. \newline {\footnotesize\color{leanink}$\checkmark$\,\code{SeedCapture.genSeqAvoid\_ne\_avoided}} \\
175 & Unrestricted generalized Mullin conjecture $\mathrm{GenMC}(n)$ over all primes {\footnotesize(S316)} {\footnotesize[revival 0]} & Define $\mathrm{GenMC}(n)$ as ``every prime appears in the generalized Euclid--Mullin sequence seeded at $n$'' and reduce MC to $\mathrm{GenMC}(2)$ via the index shift $\seq{2}(k)=\seq{}(k+1)$ (\code{gen\_mc\_two\_implies\_mc}), with cofinal walk hitting as the sufficient condition (\code{gen\_hitting\_implies\_gen\_mc\_proved}). & False for every $n\ge 2$: $n \mid \Prod{n}(k)$ for all $k$ (\code{start\_dvd\_genProd}), so a prime $p\mid n$ is coprime to $\Prod{n}(k)+1$ and is never selected. In particular $2$ never re-appears in $\seq{2}$, so the old $\mathrm{GenMC}(2)\Rightarrow\mathrm{MC}$ was $\bot\Rightarrow\mathrm{MC}$, and the cofinal-hitting premise was unsatisfiable for every $n\ge1$ (the walk is identically $0$ mod $p\mid n$; for $n=1$ the $q=2$ clause fails). Same defect as Conjecture~A (\#1). Repaired 2026-08-20: $\mathrm{GenMC}(n)$ now quantifies over primes $q\nmid n$, the hitting hypothesis likewise, and $\mathrm{GenMC}(2)\iff\mathrm{MC}$ (\code{gen\_mc\_two\_iff\_mc}). \newline {\footnotesize\color{leanink}$\checkmark$\,\code{gen\_mc\_unrestricted\_false}} \\
\midrule
\multicolumn{4}{@{}p{13.4cm}@{}}{\textsc{Technique mismatch} (TM, 55 entries) --- \emph{the tool needs structure (independence, multiplicativity, stationarity, algebraic families) that the walk provably lacks}.} \\[2pt]
4 & Consecutive vs arbitrary subsequences (ordering problem) {\footnotesize(S5)} & Cauchy--Davenport, EGZ, Davenport constant, Gordon sequenceability, Kneser, critical numbers to force partial products to cover $(\mathbb{Z}/q)^\times$ & Those theorems concern arbitrary subsequences/subset products or existence of some good ordering; DH needs prefix products of one fixed ordering (the EM walk). Generation/coverage of subsets says nothing about the specific order \\
7 & Self-avoidance implies character cancellation {\footnotesize(S6)} & Use distinctness of the multiplier primes (self-avoidance) to get cancellation in $\sum \chi(\text{walk})$ & Distinct primes all $\equiv 1 \pmod q$ give $\chi(m_k)=1$ and $|\sum P_n| = N$; self-avoidance is a prime-level property, characters see only residues; no literature exploits step distinctness \\
8 & Booker/Burgess techniques for first EM sequence {\footnotesize(S9)} & Transfer Booker's variant-EM Hasse--Weil sums, Burgess plus reciprocity for the maxFac sequence, or Booker--Simon (2026) to the minFac sequence & Booker's variant allows a choice at each step and sums over all of $\mathbb{F}_q$, not the walk; Burgess needs rapidly growing primes (largest factors); Booker--Simon proves omission for the second EM sequence only \\
13 & Additive combinatorics for EM residues {\footnotesize(S11)} & Sumset / additive-energy tools on multiplier residues to reach PED & No additive structure in EM residues; sums not products — wrong structure; Bogolyubov--Ruzsa etc. reduce to the ordering problem \\
15 & Cross-prime coupling via minimality {\footnotesize(S13)} & Use minimality ($\mathrm{seq}(n+1)$ = smallest prime factor) to couple PED failure at one modulus with walks at others, aiming for CRT contradiction & CRT independence: PED failure at $q_0$ does not constrain residues at $q_1$; minimality constrains size not residue class; multi-modulus sieve gives density-zero not finite exceptional set; all primes $\equiv 1$ mod $q$ is not self-contradictory \\
17 & Kernel-run position shift {\footnotesize(S13)} & Structural NoLongRuns route: during a kernel run the walk stays in one coset; hope revisits force a contradiction & Within a run the walk visits at most $(q-1)/d$ positions of one coset and revisits give no contradiction: products of subgroup elements stay in the subgroup \\
22 & Order-2 cancellation insufficient for MC {\footnotesize(S15)} & BRE/CCSB restricted to order-2 characters (provable from NoLongRuns by alternating-series bound; later QuadraticCharacterLargeSieve, $d=2$ Legendre) & Order-2 cancellation only separates QR from QNR; hitting $-1$ needs all non-trivial characters in the Fourier inversion, so a proper character subgroup cannot isolate the target \\
23 & Furstenberg theorem for circle rotations {\footnotesize(S15)} & Furstenberg / random-product theorems on the circle for walk equidistribution & Lyapunov exponent of unimodular scalar products is 0, so Furstenberg-type expansion gives nothing; steps are deterministic and correlated \\
24 & CCSB for bounded-order characters {\footnotesize(S15)} & Prove CCSB only for characters of bounded order and hope it suffices & Not sufficient: hit-count Fourier inversion needs every non-trivial character; low-order cancellation cannot isolate $-1$ \\
26 & PED plus self-avoidance (residue vs prime) {\footnotesize(S16)} & Combine positive escape density with the walk's self-avoidance (distinct primes) to force CCSB / a StructuredBRE & Self-avoidance constrains primes, not residue classes: $p_1 \ne p_2$, $p_1 \equiv p_2$ give identical $\chi$-values, and Dirichlet supplies infinitely many primes per class, so any character-value sequence is compatible with self-avoidance \newline {\footnotesize\color{leanink}$\checkmark$\,\code{dirichlet\_residues\_independent}} \\
27 & StructuredBRE with self-avoidance {\footnotesize(S16)} & A BlockRotationEstimate taking self-avoidance plus PED as hypotheses & Prime-to-character map is many-to-one; injectivity in the domain says nothing about the image; no version incorporating self-avoidance helps \\
29 & Counting/dispersion from self-avoidance {\footnotesize(S16)} & Count self-avoiding continuations keeping the walk clumped and argue dispersion & Same obstruction: Dirichlet provides unlimited primes in every residue class, so counting fresh primes never forces dispersion \\
31 & Large sieve for partial products {\footnotesize(S17)} & Apply additive/multiplicative large sieve to the walk character sum & Large sieve handles linear sums; the EM sum is a sum of partial products, a structure no large-sieve inequality bounds. Later: \code{ArithLSImpliesMMCSB} circular (\#61), PrimeArithLS mismatch (\#75) \\
34 & Death-set coupling across moduli {\footnotesize(S17)} & Couple death sets $\{m : \mathrm{minFac}(m) \equiv -c^{-1}\}$ at different moduli to force a hit & Death sets vary per step and reset (non-cumulative); no uniform coupling bound exists across moduli \\
39 & Linnik's theorem for MC {\footnotesize(S19)} & Use Linnik (small least prime / least QNR in APs) & Constrains prime existence, not EM factorizations; trivially $4$ is a QR mod $p \ge 5$ so QNR $\le 4$ — proved but irrelevant to walk dynamics \\
40 & Qualitative Dirichlet implies NoLongRuns {\footnotesize(S19)} & Derive NoLongRuns($L$) from Dirichlet's theorem alone & Needs density (equal divergence of $\sum 1/p$ per class), not mere infinitude; NoLongRuns requires SieveTransfer. Superseded: sieve-route endpoint GenericLPFEquidist false, SieveTransfer vacuous (\#160) \\
41 & Elementary route avoiding PNT {\footnotesize(S19)} & Reach NoLongRuns without any quantitative prime-distribution input & Impossible: NoLongRuns inherently requires quantitative prime distribution; qualitative Dirichlet gives infinitely many primes per class, not proportions; SubgroupEscape does not help kernel runs. \#160 later voided the population endpoint \\
43 & DWH to sieve-route connection {\footnotesize(S20)} & Connect combinatorial DWH (arbitrary Euclid-constrained walks) to the sieve infrastructure & DWH is about general walks / subset products, sieve route is EM-specific (minFac); independent paths; consecutive-vs-arbitrary barrier (\#4) prevents connection \\
44 & Walk telescoping for DWH {\footnotesize(S20)} & Use §37 telescoping identities to help the general DWH & Telescoping identities are EM-specific and only confirm PED alone does not give BRE for $d \ge 3$; nothing transfers to general DWH \\
47 & Diaconis--Shahshahani UBL for deterministic walk {\footnotesize(S22)} & Apply the Diaconis--Shahshahani upper-bound lemma / non-abelian mixing to bound $\|P^{*k}-U\|$ for the multiplicative walk on $(\mathbb Z/q)^\times$. & UBL bounds convolutions of probability measures with i.i.d.\ steps; the EM walk is one deterministic path with history-dependent steps. Also $(\mathbb Z/q)^\times$ is abelian, so the non-abelian machinery adds nothing. \\
48 & Furstenberg systems of multiplicative functions {\footnotesize(S22)} & Apply Furstenberg-system / Frantzikinakis--Host machinery to $n\mapsto\chi(\minFac(P(n)+1))$. & The framework needs a multiplicative function defined on all of $\mathbb Z$; the EM multiplier sequence is recursively defined and not multiplicative, so no Furstenberg system exists for it. \\
49 & Non-abelian amplification / representation theory {\footnotesize(S22)} & Pascadi-style non-abelian amplification and higher-dimensional representations to strengthen character-sum cancellation. & $(\mathbb Z/q)^\times$ is cyclic: all irreducible representations are one-dimensional characters, quasirandomness is $1$ (worst case), so there is no $SL_2$-type structure to exploit. \\
50 & CRT simultaneous equidistribution across moduli {\footnotesize(S22)} & Kowalski et al.\ style joint equidistribution over many moduli via CRT to obtain multi-modular control of the walk. & Requires independence across moduli; the single integer $\minFac(P(n)+1)$ drives every residue walk simultaneously, so the EM walk couples all moduli and the CRT-independence hypothesis fails. \\
51 & Polynomially-defined multiplicative function theory {\footnotesize(S22)} & Singha Roy (2024) equidistribution results for multiplicative $f$ with $f(p)=F(p)$ polynomial, applied to the least prime factor. & $\minFac$ is not multiplicative ($\minFac(210)\ne\minFac(6)\minFac(35)$), and those results only reach moduli up to $(\log x)^K$; wrong function class. \\
53 & LSD method for recursive minFac sequences {\footnotesize(S22)} & Landau--Selberg--Delange mean-value method via Dirichlet series / $L$-functions applied to the recursive $\minFac$ sequence. & LSD handles multiplicative functions through Dirichlet series; the recursive $\minFac(P(n)+1)$ sequence is neither multiplicative nor an $L$-function coefficient sequence, so no analytic continuation input exists. \\
55 & PNT in APs formalization blocked {\footnotesize(S27)} & Formalize quantitative PNT in arithmetic progressions (\code{IK.DirichletPNT}) as input to SieveEquidistribution / BV routes. & Zero-free region $\sigma\ge 1-c/\log(q|t|)$ absent from the entire Lean 4 ecosystem; needs $\sim$5000+ lines including Hadamard factorization for $L$-functions. Even if available it is a population statement (\#90). \\
56 & Siegel--Walfisz formalization blocked {\footnotesize(S27)} & Formalize Siegel--Walfisz (\code{IK.SiegelWalfisz}) for uniform error terms in APs. & Same blocker as \#55 (zero-free region and Siegel's theorem absent from Lean 4); and it yields SieveEquidistribution (population), never orbit-level SieveTransfer. \\
57 & Direct Bombieri--Vinogradov formalization {\footnotesize(S27)} & Prove Bombieri--Vinogradov in Lean to feed the chain BV $\to$ MMCSB $\to$ MC. & Estimated 5000--10000 lines (multi-year, Vaughan identity and large sieve absent), and even a formal BV leaves the generic-to-EM transfer \code{BVImpliesMMCSB} (SieveTransfer) unsolved. \\
59 & Davenport-constant zero-sum subset products {\footnotesize(S37)} & Use $D(\mathbb Z/(q-1))=q-1$ (any $q-1$ elements contain a zero-sum subsequence) to force products of EM multipliers to reach $-1$. & DynamicalHitting needs PREFIX (consecutive) products reaching $-1$, not existence of some zero-sum subsequence; same ordering obstruction as \#4. Only marginal information, DH needs joint (position, multiplier) data. \\
60 & Kneser sumset growth for prefix products {\footnotesize(S37)} & Kneser's theorem $|A+B|\ge|A|+|B|-|H|$ to force growth of the reachable product set of the walk. & Gives weaker bounds than direct arguments and faces the same ordering obstruction (\#4) and generation-vs-coverage (\#130); Kneser is not in Mathlib (LeanCamCombi only). \\
62 & Full ALS by refining WeakALS {\footnotesize(S36)} & Interpolate from the proved WeakALS constant $N(\delta^{-1}+1)$ to the optimal large-sieve constant $N-1+\delta^{-1}$. & WeakALS treats evaluation points independently; the optimal constant exploits interference (Selberg--Beurling majorants, duality + Hilbert). No smooth interpolation exists between the two arguments. \\
63 & KernelRowSumBound from geometric-sum tools {\footnotesize(S36)} & Prove the sharp kernel row-sum bound (Schur test input for optimal ALS) using only Jordan's inequality and geometric-sum bounds. & Insufficient: needs Selberg--Montgomery--Vaughan extremal functions ($\sim$300--500 lines, not in Mathlib); S60 notes sharp KRSB is harder than ALS itself. WeakALS already suffices downstream. \\
74 & Mathlib BirkhoffSum / Birkhoff ergodic theorem {\footnotesize(S38)} & Use \code{Mathlib.Dynamics.BirkhoffSum} time-average API for the walk's visit averages. & The API assumes an orbit under a single map $T$; the EM walk applies a different multiplier at each step (non-autonomous), so no single transformation, invariant measure or Birkhoff average applies. \\
75 & PrimeArithLS to MMCSB structural mismatch {\footnotesize(S39)} & With PrimeArithLS proved, choose coefficients ($a(n)=1$ or walk-adapted $a(n)=\chi(w(n)n^{-1})$) to extract walk character-sum bounds. & PrimeArithLS sums over integer indices $\sum_{n\le N}a(n)\chi(n)$; MMCSB needs sums at walk positions $\sum\chi(w(n))$, and $w$ is many-to-one, not invertible in $(\mathbb Z/p)^\times$. Every coefficient choice returns the trivial $N$ bound. \\
78 & Partial HOD (even h=2) unprovable {\footnotesize(S40)} & Prove HigherOrderDecorrelation for small $h$ (say $h=2$) from known results on correlations of the least prime factor $P^-(n)$. & Requires the joint distribution of consecutive least prime factors $(P^-(n),P^-(n+h))$, an open problem; Tao--Ter\"av\"ainen, Pilatte, Charamaras--Richter cover additive or completely multiplicative functions only, and $P^-$ is neither. \\
86 & Nonstationary ergodic theorems (Monakov, Ito--Kawada) {\footnotesize(S51)} {\footnotesize[revival 2]} & Monakov (2024) Ces\`aro convergence of nonstationary random walks on compact abelian groups to Haar under strict aperiodicity, applied to the walk on $(\mathbb Z/q)^\times$. & Requires strictly aperiodic probability measures AND independent steps: for a Dirac mass $\delta_a$, $|E\chi(X_n)|=|\chi(a)|=1$, so strict aperiodicity fails; and $m(n)$ depends on the whole history. Both conditions fail independently. \\
89 & Michelen--Sahasrabudhe anti-concentration / PGF zeros {\footnotesize(S57)} & MS-style anti-concentration via zero-free regions of the probability generating function of the walk's empirical measure (last unexplored technique). & Anti-concentration is a lower bound on variance, the wrong direction for equidistribution; the PGF $f(z)=N^{-1}\sum z^{L(n)}$ has zeros encoding the trajectory (circular); no factorization without independent steps. Introduces the Four-Way Blocker. \\
94 & Fiber uniformity from SE alone {\footnotesize(S61)} & Derive uniform visit counts $V_N(c)=\Theta(N/(q-1))$ from the proved SubgroupEscape. & SE gives eventual coverage, not visit statistics; uniformity needs mixing-time estimates (unavailable for deterministic walks) or joint $(w(n),m(n))$ equidistribution, which is CME. Multipliers $\{2,3\}$ alternating in $(\mathbb Z/5)^\times$ generate yet avoid $-1$ forever. \\
95 & Spectral gap for deterministic walks {\footnotesize(S62)} {\footnotesize[revival 0]} & PED gives step-distribution spectral gap $\rho<1$ (proved abstractly); infer CCSB, or apply Markov mixing/spectral theory to the EM walk. & Spectral gap constrains frequency of step types, not their ordering: clumped kernel-then-escape steps give walk sum $\Theta(N)$. Spectral theory applies to distributions (random sampling), not a single deterministic path; each EM transition operator is a permutation matrix with unimodular spectrum. \\
109 & Non-multiplicative Halasz {\footnotesize(S81)} {\footnotesize[revival 0]} & Extend Halasz / pretentious-distance theory to minFac-driven partial products. & No extension exists (2024--2026 searches); pretentious distance is intrinsically Euler-product based and minFac is provably not multiplicative (Four-Way Blocker item 2). Tao--Teravainen correlations stay within the multiplicative world. \\
111 & Rough-number concentration for $d=2$ {\footnotesize(S82)} & For the quadratic character use coprimality, $q$-roughness and super-exponential growth of Euclid numbers to forbid $L$ consecutive QR-valued minFacs (NoLongRuns). & For any $L$ one builds $L$ pairwise coprime $q$-rough integers of arbitrary size all with QR minFac (distinct QR primes times large coprime cofactors). Only the recursion is EM-specific and using it to control minFac residues is SieveTransfer; coprimality does not imply minFac independence. \\
118 & Growth gives no decorrelation {\footnotesize(S137)} & DSL variance route: derive quantitative decorrelation of population cross-terms $C(j,k,X)$ from super-exponential growth of the accumulator. & Residues mod $q$ are periodic regardless of magnitude, so growth is invisible mod $q$; the $+1$ shift is arithmetically entangled; CRT invariance is structural, not statistical. Every dynamical mechanism reduces to sieve content (JSE via CRTPropagationStep). Superseded: variance hypotheses false (\#156), DSL vacuous (\#160). \\
126 & UCE is PE restated in progression language {\footnotesize(S165)} & Uniform Conductor Equidistribution: $\minFac N \bmod q$ equidistributed among $N \equiv 1 \pmod M$ uniformly in the conductor $M$, via $P(n)+1 \equiv 1 \pmod{\prod n}$. & At $X \approx M_n$ the counting set is the singleton $\{P(n)+1\}$; density over a singleton is the pointwise MC problem, so \code{UCEImpliesCME} is the population-to-orbit gap (\#90) and $X/M = O(1)$ is outside every sieve asymptotic (\#108). Superseded: UCE is false (\#160). \\
128 & p-adic geometry, diamonds, Hecke orbits inapplicable {\footnotesize(S167)} {\footnotesize[revival 1]} & Perfectoid spaces, Scholze--Weinstein diamonds, Fargues--Fontaine slopes, and Hecke-orbit equidistribution (Clozel--Ullmo, Eskin--Mozes--Shah) applied to the EM walk. & Every geometric equidistribution theorem needs a fixed algebraic correspondence or group action; the EM multiplier is state-dependent, history-dependent and non-algebraic. The FF curve lives at a single prime while the walk is adelic; the slope analogy is a category error. \\
131 & Dobrushin coefficient equals one {\footnotesize(S172)} {\footnotesize[revival 0]} & Relax CME to \code{MultiplierUniformityBound} and use Dobrushin non-homogeneous Markov convergence plus a stopping-time reading of $\minFac$. & Each EM transition kernel is a point mass, so $\alpha(K_n)=1$ and $\sum(1-\alpha_n)=0$: Dobrushin's theorem never applies. Batching stays deterministic, windowing is CME again, conditioning reintroduces selection bias; MUB is weaker than CME but vacuous; stopping times repackage DH. \\
133 & Self-similar functional equation mismatch {\footnotesize(S173)} {\footnotesize[revival 0]} & Use the tail identity (orbit from $\prod M$ = tail of standard orbit) as self-similarity in the Lapidus--van Frankenhuijsen fractal-zeta framework $\zeta(s)=\sum r_j^s\zeta(s)+f(s)$. & The identity gives $L_{\mathrm{EM}} = \text{head} + L_{\text{from }\prod M}$, but the tail orbit starts at $\prod M \neq 2$ and is a different series, not a scalar multiple of $L_{\mathrm{EM}}$; there are no scaling ratios, so the framework does not apply. \\
135 & Number field extensions invisible to walk {\footnotesize(S179)} {\footnotesize[revival 1]} & EM over $\mathbb{Z}[i]$ (norm-1 subgroup of $\mathbb{F}_{p^2}^\times$), biquadratic fields, Hecke Grössencharacters, \code{NormTwistedCME}, arbitrary $\mathcal{O}_K$. & Universal Confinement: $\mathbb{Z} \to \mathcal{O}_K/\mathfrak{p}$ factors through the prime subfield $\mathbb{F}_r$, so every character of $(\mathcal{O}_K/\mathfrak{p})^\times$ restricts to a Dirichlet character mod $r$ on the integer walk; Hecke characters add only archimedean growth; inert primes give more characters, not fewer. \\
141 & Rogue character contradicts no known bound {\footnotesize(S283)} {\footnotesize[revival 0]} & Rigidity framework: perpetual avoidance of $q$ yields $\chi\neq 1$ with $\|S_\chi(N)\| \ge N/(q-2)$; contradict it with a character-sum estimate. & No standard bound applies to $\sum_{n<N}\chi(\prod n \bmod q)$: Pólya--Vinogradov (not consecutive integers), Burgess (no interval), Halász (not multiplicative in $n$), large sieve (averages over $\chi$), Shiu (consecutive primes in one class exist), CRT (moduli independent); a walk-specific bound is CCSB itself. \newline {\footnotesize\color{leanink}$\checkmark$\,\code{perpetual\_avoidance\_rogue}} \\
144 & Reciprocity transfer to min sequence fails {\footnotesize(S298)} {\footnotesize[revival 0]} & Transport the Cox--van der Poorten / Booker omission proofs (Jacobi symbols between sequence data, moduli growing with the orbit) from max to min sequence. & Max proofs make a real character constant on the finite factor support of the Euclid number; $\minFac$ confines it to a cofinite set, where no nontrivial character is constant (\code{char\_non\_constancy}). The invariant is congruential at $\Pi_n=8mP_n$: eviction automatic, fullness at $\Pi_n$, fragment empty. \newline {\footnotesize\color{leanink}$\checkmark$\,\code{Reciprocity.no\_reciprocity\_induction\_proof}} \\
147 & Avoidance does not force factor-set diversity {\footnotesize(S299)} {\footnotesize[revival 1]} & Contrapositive of the diversity chain: a missing $q$ forces eventually monochromatic factor sets mod $q$, contradicting \code{diverse\_steps\_imply\_vanishing}. & (F1) \code{meanCharValue} contracts by averaging over a factor set while the walk selects one factor with $\|\chi(s)\|=1$; (F2) the chain concludes some branch reaches $-1$, avoidance constrains one branch; (F3) \code{productMultiset} fixes factor sets in advance, real ones are path-dependent. \newline {\footnotesize\color{leanink}$\checkmark$\,\code{diverse\_steps\_imply\_vanishing}} \\
152 & Support-invisibility of algebraic invariants {\footnotesize(S299)} {\footnotesize[revival 0]} & Find an algebraic invariant of the collected prime bag that detects which primes are missing. & Missingness is a support condition on $\sum_p e_p$, but every computable algebraic invariant factors through $p\mapsto p \bmod m$, hence through the walk, which sees only the product. \\
167 & Logarithmic density for the simultaneous-in-$q$ form {\footnotesize(S312)} {\footnotesize[revival 0]} & Reformulate the seed-average law in logarithmic density, on the grounds that logarithmic density is compatible with countable operations in ways natural density is not, then close the simultaneous-in-$q$ form by Borel--Cantelli. & The premise is false: logarithmic density is \emph{also} only finitely additive (singletons have log density $0$ and their union is $\mathbb{N}$), so it supports no Borel--Cantelli. The settings where log density genuinely is better behaved do not apply: Davenport--Erd\H{o}s needs \emph{sets of multiples} (the failure sets $F_q$ are not divisibility conditions --- that is exactly what excluding $q$ from \code{bandUpTo} means, \#165), and log-averaged Furstenberg systems gain from \emph{dilation invariance}, which the greedy map lacks and which is orthogonal to simultaneity anyway. Independently, the $1/m$ weights destroy \code{SelectionLaw.selection\_law}: it is an \emph{equality} obtained from CRT surjectivity on a full period and requires a weight constant on residue classes, whereas $\sum_{m\le X,\ m\equiv a\,(M)} 1/m = (1/M)\log(X/M)+O(1/M)$ is exact only as $X/M\to\infty$ --- precisely the regime in which log density coincides with natural density on these $M$-periodic events. Pure loss. \newline {\footnotesize\color{leanink}$\checkmark$\,\code{SelectionLaw.selection\_law}} \\
168 & Summable per-$q$ rates close the union over all primes {\footnotesize(S312)} {\footnotesize[revival 1]} & Make the Theorem C failure fraction summable in $q$ by choosing $n(q)$, then union-bound (Borel--Cantelli) over all primes to get the simultaneous-in-$q$ form of the seed-average law. & Summability \emph{is} achievable ($\varepsilon_q=q^{-2}$ with $C_c(q)=\max(48q,3q^2)$ and $n(q)\approx 2e^{250}\exp(8B\varphi(q))$; the $e^{25}\log n/n$ tail is slack, not binding), and the differing sample spaces are \emph{not} an obstruction ($\mathrm{modulus}\,q\,Y \mid P(Y')$, and \code{SelectionLaw.genSeqAvoid\_prefix\_eq\_of\_modEq} is stated for an arbitrary multiple of the band primes), so the finite-$S$ union bound is a genuine corollary (\code{AlmostAllDensity.finite\_simultaneous\_density}). But at scale $X$ only $q \lesssim \log\log\log X$ are controlled, and Borel--Cantelli is a theorem about \emph{countably additive} measures: under a finitely additive density $\sum_q \varepsilon_q<\infty$ implies nothing about $\bigcup_q F_q$ (increasing sets of density $\le\delta$ can have union of density $1$). The rate is not the obstruction; the additivity is. The repair is a countably additive ambient measure --- the profinite ensemble, in which every event is a cylinder event; built in Session 314 as $\Omega=\prod_r \mathbb{Z}/r$ with the product uniform measure (squarefree moduli make $\mathbb{Z}_r$ unnecessary), where the union over all $q$ is a one-line consequence of countable additivity (\code{ProfiniteHeadline.measure\_some\_prime\_missed\_eq\_zero}). That is a statement about a different ensemble, in which $\mathbb{N}$ is null; it does not revive the natural-density form. \newline {\footnotesize\color{leanink}$\checkmark$\,\code{AlmostAllDensity.finite\_simultaneous\_density}} \\
174 & Bounding the missed-prime counting function {\footnotesize(S313)} {\footnotesize[revival 1]} & Bound $\mathrm{missed}(x)=\pi(x)-\#\{k: p_k \le x\}$, or $\sum_{q \text{ missed} \le x} 1/q$, or any nontrivial function of the missed set of one orbit. & Every such target requires a \emph{lower} bound on the hit count. The sure layer's only multiplier-size statement is \code{few\_small\_multipliers}, an \emph{upper} bound on the number of small multipliers --- the wrong direction, and the only one available, because distinctness is the only sure arithmetic fact about the multiplier sequence. Moreover a lower bound is gated on an open problem: under \code{AutonomousBranch.PerpetualPrimality} (open; its negation is (C$\infty$), the top frontier item) $p_k = \prod_k+1$ grows doubly exponentially and $\mathrm{hits}(x)=O(\log\log x)$; even granting (C$\infty$), compositeness gives only $p_k \le \sqrt{\prod_k+1}$, still doubly exponential. Hence any nontrivial bound on $\mathrm{missed}(x)$ implies (C$\infty$) and much more, and cannot come from a bookkeeping layer. Underlying cause: the \emph{sign asymmetry of} \code{minFac} --- $p_k=\minFac(N_k)$ yields infinitely many non-divisibility facts and exactly one divisibility fact, about a prime captured by definition. \\
177 & Height argument for capture at a fixed prime {\footnotesize(S318)} {\footnotesize[revival 1]} & Build a functional on the greedy orbit from absolute values $|\cdot|_v$, $q$-adic logarithms and the product formula (the greedy step is multiplication by the principal idele $p_{n+1}$) such that the failure of MC at $q$ --- $E_n\in\mathbb Z_q^\times$ eventually (\code{AdelicShadow.misses\_iff\_eventually\_unit\_both}) --- violates a size bound; via Ridout/Roth exponents, the $S$-unit equation $E_n-P_n=1$ with Baker/subspace, or the product of Euclid numbers in $\Omega$. & The approximation exponent $\prod_{v<\infty}|P_n+1|_v=1/E_n$ is identically $H^{-1}$ (below Roth's $2$, no information); Baker degrades as $\exp(-C^{\omega}\prod h)$ with $\omega\approx n$ terms, far below $1/E_n$, and the subspace theorem needs a fixed support while $P_n$ is an $S_n$-unit with $S_n$ growing (\#134 mismatch); $\prod E_n\to0$ in $\Omega$ is the tautology ``every prime divides some Euclid number''. The structural reason: $\mathbb Z_q^\times=\mu_{q-1}\times(1+q\mathbb Z_q)$, the walk mod $q$ is the \emph{torsion} (Teichm\"uller) component, and every absolute value and $q$-adic logarithm factors through the torsion-free data. Formally, the capture identity shows that for a fixed multiplier prefix --- hence fixed heights, defects and all $|\cdot|_v$, $v\ne q$ --- capture of $q$ is a condition on the single residue $m\bmod q$ with both outcomes realised, so no inequality among height-type functionals decides it. The only height/torsion coupling left (Northcott: bounded height $\Rightarrow$ finitely many) reproduces the seed-average programme on the bounded-height family $m\le X$. See \code{docs/analysis/height\_argument\_attempt\_2026-08-20.md}. \newline {\footnotesize\color{leanink}$\checkmark$\,\code{SeedCapture.captured\_iff\_mem\_visited}} \\
\midrule
\multicolumn{4}{@{}p{13.4cm}@{}}{\textsc{Scale mismatch} (SM, 14 entries) --- \emph{the error term of the tool exceeds the signal on an $O(\log x)$-term or exponentially sparse orbit}.} \\[2pt]
11 & BlockDecorrelation for high-order characters {\footnotesize(S11)} & Block-averaged decorrelation of multiplier character values as intermediate toward PED/CCSB & Works for low-order characters but fails for order $d \sim q-1$: block bound does not scale with the character order \\
28 & Abel summation for walk sums {\footnotesize(S16)} & Summation by parts to convert multiplier decorrelation into walk-sum bounds & Linear weights amplify contributions: yields $O(N^2)$ not $o(N)$, strictly worse than trivial \\
33 & CRT product-group reformulation {\footnotesize(S17)} & Run the walk on $\prod_{q \le Q} (\mathbb{Z}/q)^\times$ simultaneously & Makes the problem harder: the product group is exponentially large; no gain over single modulus. Later \#101 (profinite adds nothing) \\
38 & Dec to CCSB via Abel/VdC $h=1$ {\footnotesize(S18)} & Convert multiplier decorrelation to a walk-sum bound by Abel summation or van der Corput with the single shift $h=1$ & Product-to-sum conversion fails; VdC with $H=1$ gives $|S_N|^2 \le N^2/2 + o(N^2)$, i.e. $|S_N| \le N/\sqrt2$, a constant fraction not $o(N)$; $h \ge 2$ needs HOD \\
77 & Dec plus VdC insufficient for CCSB {\footnotesize(S40)} & Van der Corput with only the $h=1$ autocorrelation, i.e.\ from DecorrelationHypothesis alone. & VdC at $H=1$ gives $\|S_{\mathrm{walk}}\|\le N/\sqrt2$, a constant fraction of the trivial bound, not $o(N)$; $o(N)$ needs autocorrelation control for all lags $h\le H$ with $H\to\infty$. \\
80 & Order-2 CCSB from PED plus NoLongRuns {\footnotesize(S41)} & For quadratic characters exploit the alternating kernel-block structure: PED gives linearly many escapes, NoLongRuns bounds block length $L$; hope the alternating block sum is $o(N)$. & $|S|\le(\text{number of pairs})\times L=O(N)$, not $o(N)$; adversarial block lengths $L_1=2,L_2=0,L_3=2,\dots$ give alternating sum $N$ with $\Omega(N)$ blocks. Alternation alone forces no cancellation below $O(N)$. \\
91 & Prime Euclid density argument {\footnotesize(S59)} & Use the sub-population of prime Euclid numbers $P(n)+1$ (multiplier equals the whole number) as leverage on the walk or CME. & Under doubly-exponential growth ($\log P_n \asymp 2^n$, the perpetual-primality branch) $\sum 1/\log P_n$ converges and prime Euclid numbers are heuristically finitely many; on the branch MC predicts ($\log P_n \asymp p_n$) the heuristic count diverges (paper \S4.1). Either way they contribute $o(N)$ and say nothing about the composite population, which is the irreducible content. \\
96 & Level-of-distribution scale mismatch {\footnotesize(S63)} {\footnotesize[revival 0]} & Prove a Bombieri--Vinogradov type level of distribution for EM products (EMHasLevelOfDistribution) and transfer it to MMCSB walk character sums. & LoD error $(\mathrm{prod}\,N)^\theta/(\log\mathrm{prod}\,N)^A\ge 2^{\theta N}$ grows exponentially while MMCSB needs $\varepsilon N$; \code{LoDImpliesMMCSB} is false. Any walk-adapted LoD reduces to CCSB (walk sums) or Dec (multiplier sums); no intermediate. \newline {\footnotesize\color{leanink}$\checkmark$\,\code{prod\_superlinear}} \\
102 & Large sieve single-frequency extraction {\footnotesize(S72)} & Apply the analytic large sieve at the single frequency $\alpha=0$ ($R=1$) to bound the walk sum $S_N$. & With one evaluation point the separation condition is vacuous, giving the trivial $|S_N|\le N$; embedding $\alpha=0$ in $R+1$ points gives $|S(0)|^2\le N^2(R+1)$, worse than trivial. Large sieve averages over frequencies; it cannot extract pointwise information on one walk. \\
108 & Harper BDH / Weil scale mismatch {\footnotesize(S81)} {\footnotesize[revival 1]} & Apply Harper's Brownian-Dirichlet variance asymptotics for non-multiplicative sequences to EM products; registry restates over $\mathbb F_p[t]$ via the Weil bound. & BDH needs well-distribution in APs (circular for EM) and non-concentration; EM products are super-exponentially sparse, and it yields variance over most $q$, not pointwise. Weil error $O(p^{n/2})$ per degree is vacuous for an orbit at a single degree per step. \\
154 & LSD/Wirsing density along the orbit {\footnotesize(S299)} {\footnotesize[revival 1]} & Landau--Selberg--Delange / Wirsing mean-value theorems for a density statement along the EM orbit. & No exponentially-sparse LSD theorem exists; the orbit contributes $O(\log x)$ terms below $x$, far under every LSD error term (Wirsing 1961, Tenenbaum II.5, Serre 1976 checked: nothing for sparse sequences). \\
155 & Nonstandard / ultraproduct receptacle {\footnotesize(S299)} {\footnotesize[revival 0]} & Transfer the avoidance statement to a hyperfinite orbit by Łoś and extract a Loeb-measure Gap. & Detection is honest (Łoś transfers avoidance) but the Loeb measure of the hyperfinite orbit is $0$ for every sequence, avoiding or not; a conservative extension yields no new Gap by definition. \\
161 & Seed-magnitude large-step threshold {\footnotesize(S309)} {\footnotesize[revival 3]} & In the seed-average programme, define the large-step threshold from the accumulator size, $y_k = C k \log_2 P_k$, so that a ``large step'' at depth $k$ means $\tilde p_k > y_k$. & $y_k$ is a function of the seed's magnitude, not of its type, so $\{\tilde p_k > y_k\}$ is not measurable for the type $\sigma$-algebra mod $M_Y$; forcing measurability by size-stratification needs $Y \gtrsim (\log X)^A$, i.e.\ modulus $\exp((\log X)^A) \gg X$ --- the Bombieri--Vinogradov regime (\#96). Independently, the conditional probability of a large step at that threshold is $\asymp 1/\log\log X$, so Theorem~C would be vacuous and (LS) false as stated. Repaired by the type-determined $y_k = C k \log_2 c_k$ (cofactor), which keeps the modulus constant and gives $\log y_k = (2+o(1))\log k$, exactly what the far-band estimate needs. \\
163 & Truncation quantifier order in Theorem C {\footnotesize(S309)} {\footnotesize[revival 3]} & State the truncation limit of Theorem~C as $\exists Y_0\ \forall Y \ge Y_0$ at a fixed horizon $n$ (the (e-3) shape, and \S F.2 of the (LS) proof). & Under that order the far-band constant $c_0$ degenerates as $\log\log Y/\log k \to \infty$, so $c_0$ is not absolute; and simultaneously forcing the degenerate-prefix tail to vanish requires $\log Y \gg n \log n$. The two requirements are incompatible at fixed $n$. Repaired by tying $Y$ to $n$ ($\log Y \asymp n^2$, cutoff $k \ge n/\log n$), which gives $c_0 \ge e^{-23}$ and changes the statement shape of Theorem~C. \\
\midrule
\multicolumn{4}{@{}p{13.4cm}@{}}{\textsc{Methodological rules} (MR, 3 entries) --- \emph{not mathematical obstructions but standing rules of the project (no numerical certificates, no instance-by-instance work), numbered in the earliest sessions}.} \\[2pt]
1 & Computing concrete sequence values {\footnotesize(S1)} & Compute $\mathrm{seq}(8)$--$\mathrm{seq}(12)$ or other concrete Euclid--Mullin terms inside the Lean development & Methodological rule: concrete values prove nothing about all primes $q$; breaks the build, contributes nothing (Guiding Principle). Closest code OS: finite data says nothing about the whole orbit \\
2 & Raising ThresholdHitting by computation {\footnotesize(S1)} & Increase $T$ in \code{ThresholdHitting T} $\Rightarrow$ MC by checking individual primes & Checking individual primes does not prove the conjecture; $T=11$ is the proved reduction (\code{threshold\_11\_implies\_mullin'}) and no finite $T$ closes MC \\
3 & Brute-force SubgroupEscape instances {\footnotesize(S1)} & Add \code{SubgroupEscape} instances prime by prime beyond $q \le 157$ & SE is already proved globally via PRE $\Rightarrow$ SE; instance-by-instance work is redundant with the global theorem \\
\midrule
\multicolumn{4}{@{}p{13.4cm}@{}}{\textsc{Unassigned numbers} (10): \emph{numbers that were never assigned to an entry in any session log (kept so the numbering is accounted for)}: \#25, 64, 65, 66, 67, 68, 69, 70, 71, 72.} \\
\end{longtable}
\renewcommand{\arraystretch}{1.0}\normalsize
